我们需要如下近单调性，它在大筛法递推中也被用到，见 \cite[引理~4.4 和注记~5]{GL} 与
\cite[引理~3.1]{BGL}：
\begin{equation}\label{eq:sextic-sieve-norm-monotone}
 D_j(M_1,N)\ll_S D_j(M_2,N)
 \quad\text{if }M_2\ge C M_1\log(2M_1N),\quad M_1,M_2,N\ge1.
\end{equation}
下面给出一个直接验证。对于零范数，该断言是显然的。
否则，为 $D_j(M_1,N)$ 选取一个使系数取极大的单位系数向量，
并在 $3M_1/2$ 处把行壳层拆分；其中一部分
至少承载其一半的质量。令 $L=M_2/M_1$。对第一部分使用范数落在
$(L,4L/3]$ 中的好素理想，对第二部分使用范数落在
$(2L/3,L]$ 中的素理想。于是当 $p\nmid a$ 时 $ap\in\mathcal A_{M_2}$。固定域上的
素理想定理给出 $P\gg_S L/\log L$ 个这样的素数。范数不超过 $2x$ 的理想至多包含其中
$d_x=\log(2x)/\log(2L/3)$ 个。对
$S_a(v)=\sum_bv_bF_b(a)^j$ 与 $p\nmid a$，乘性给出
\[
 S_a(\lambda)=S_{ap}\bigl((\lambda_b\overline{F_b(p)^j})_b\bigr)
                  +S_a\bigl((\lambda_b1_{p\mid b})_b\bigr).
\]
第一个系数在 $p\mid b$ 处为零，否则消去 $p$ 处的因子。对所选取的行与
素数取平方并求和，得到
\[
 \frac{P-d_{M_1}}2D_j(M_1,N)
 \le2P D_j(M_2,N)+2d_ND_j(M_1,N).
\]
对所述的阈值，当 $C$ 充分大时
$d_{M_1}+4d_N\le P/2$，这就证明了
方程~\eqref{eq:sextic-sieve-norm-monotone}。

我们以 \cite[引理~3.4--3.5]{BGL} 的幂分解与递推作结。假设对 $j\in J$ 同时地，
对每个正的损失
\begin{equation}\label{eq:sextic-sieve-inductive-bound}
 D_j(M,N)\ll_{S,\varepsilon}(MN)^\varepsilon
       \{M^\alpha+N+(MN)^{2/3}\},\qquad \alpha>4/3,
\end{equation}
对 $M,N\ge1$ 成立；单位壳层直接满足同样的界。初始估计以 $\alpha=2$ 给出这一情形。
把好的赋值理想唯一地写成
\[
 a=a_1a_2^2a_3^3a_4^4a_5^5a_6^6,
\]
其中 $a_1,\ldots,a_5$ 两两互素且无平方因子，$a_6$ 任意。代入
$a_i\in\mathcal I_{X_i}$ 并令 $X=\prod_iX_i$；这里 $X_i\ge1/2$ 且
$\prod_iX_i^i\asymp M$。选取 $\ell\in\{1,2\}$ 使得
$X_\ell=\max(X_1,X_2)$。在其余各因子固定后，系数
\[
 \lambda'_b=\lambda_b\prod_{i\ne\ell}F_b(a_i)^{ij}
\]
的绝对值至多为 $|\lambda_b|$，而剩余的核恰为 $D_{j\ell\bmod6}(X_\ell,N)$ 的核。
放大其固定的正行限制便得到这个范数界。这包含了六次幂因子的零点。
集合 $J$ 在映射 $j\mapsto j\ell\bmod6$ 下封闭。
对固定的理想计数，并在正和式中忽略它们的互素性，得到
\begin{equation}\label{eq:sextic-sieve-norm-power}
 E_j(M,N)\ll_{S,\varepsilon}(MN)^\varepsilon
 \max_{\boldsymbol X}
 \min\left\{M^{1/3},(M/X_\ell)^{2/3}\right\}
             D_{j\ell\bmod6}(X_\ell,N).
\end{equation}
为验证该因子，固定选择的数目是 $O(X/X_\ell)$。
当 $\ell=1$ 时，它至多为一个常数倍的同时满足
$(M/X_1)^{1/2}\ll(M/X_1)^{2/3}$ 与
$(MX_2/X_1)^{1/3}\le M^{1/3}$ 的量；当 $\ell=2$ 时，它至多为一个常数
倍的 $(MX_1^2/X_2^2)^{1/3}\le M^{1/3}\ll(M/X_2)^{2/3}$。
盒子的数目是 $\log(2+M)$ 的一个固定幂。

把方程~\eqref{eq:sextic-sieve-inductive-bound} 代入
方程~\eqref{eq:sextic-sieve-norm-power}，得到
\[
 E_j(M,N)\ll_{S,\varepsilon}(MN)^\varepsilon
                    (M^\alpha+M^{1/3}N).
\]
化简前的三项分别被 $M^\alpha$、$M^{1/3}N$ 与 $(MN)^{2/3}$ 界住；
最后一项依 $N\le M$ 或 $N\ge M$ 分别被前两项之一界住。把这一估计用于
方程~\eqref{eq:sextic-sieve-norm-reflection}。所有辅助的非空
尺度都被 $2+MN$ 的固定幂界住：这里 $N/\Delta_{\mathrm d}\ge1/2$ 且
$H(N/\Delta_{\mathrm d})\ll_{S,\varepsilon}(2+MN)^{3+\varepsilon}$。
单位列壳层直接满足同样的估计。因此
可以把预备损失取得更小，以给出任意指定的最终损失。
两个尺度因子为
\[
 \Delta_{\mathrm r}^{\alpha-1}\Delta_{\mathrm d}^{1-2\alpha}
       \ll_S\Delta_{\mathrm d}^{-\alpha}\le1,
 \qquad (\Delta_{\mathrm r}\Delta_{\mathrm d})^{-2/3}\le1.
\]
因此，对 $M,N\ge1$，
\begin{equation}\label{eq:sextic-sieve-norm-recurrence}
 D_j(M,N)\le E_j(M,N)\ll_{S,\varepsilon}(MN)^\varepsilon
       \{M+N^{2\alpha-1}M^{1-\alpha}+(MN)^{2/3}\}.
\end{equation}

令 $h=2-3/(3\alpha-1)=(2\alpha-5/3)/(\alpha-1/3)$ 且
$\beta=2-2/(3\alpha-1)$。若 $M\ge N^h$，
则方程~\eqref{eq:sextic-sieve-norm-recurrence} 中的中间项
至多为 $(MN)^{2/3}$。若 $M<N^h$，则把
方程~\eqref{eq:sextic-sieve-norm-monotone} 中的 $M$ 换成
$N^h\log(2MN)$ 的一个固定倍数后应用之。那里的方程~\eqref{eq:sextic-sieve-norm-recurrence}
给出 $(MN)^\varepsilon N^\beta$，因为
$2(h+1)/3=\beta$ 且 $h<\beta$。于是两种情形下都有
\[
 D_j(M,N)\ll_{S,\varepsilon}(MN)^\varepsilon
                   \{M+N^\beta+(MN)^{2/3}\}.
\]
对称性以 $\beta$ 代替 $\alpha$ 给出方程~\eqref{eq:sextic-sieve-inductive-bound}，
且对每个 $j\in J$ 同时成立。
从 $2$ 出发，映射 $\alpha\mapsto2-2/(3\alpha-1)$ 递减趋于
$4/3$。对于所要求的损失，取有限多次迭代，直到
剩余的指数间隙小于该损失的一个固定比例，再
把它吸收进 $(MN)^\varepsilon$。结合该界与其对称
形式，在 $M\le N$ 时利用
$M^{4/3}\le(MN)^{2/3}$，在 $N\le M$ 时利用相应的不等式，
即证得
\begin{equation}\label{eq:sextic-sieve-norm-shell}
 D_j(M,N)\ll_{S,\varepsilon}(MN)^\varepsilon
               \{M+N+(MN)^{2/3}\}\qquad(j\in J).
\end{equation}

对两个范数球作二进细分，并沿列
壳层使用 Cauchy 不等式，可把方程~\eqref{eq:sextic-sieve-norm-shell} 转化为元素为
$F_b(a)=\chi_a(b)$ 的球矩阵的同样
界。对数
因子被一个更小的预备损失吸收，而单位壳层
则直接界住。方程~\eqref{eq:sextic-large-sieve}
中 $\varsigma=1$ 的矩阵是这一矩阵在长度 $(D,K)$ 处的转置。
复矩阵与其转置具有相同的算子范数，这由
伴随恒等式继之以全向量共轭得到。该界关于两个长度
对称，因此它恰好给出
$K+D+(KD)^{2/3}$。对整个内和取共轭得到
$\varsigma=-1$，零点不变。最后，固定的行限制
只会去掉非负项，而固定的系数限制通过
零延拓来实现。这就证明了引理的全部断言。
\end{proof}

\subsection{物理行及其逆见证}

大筛法既已证得，剩下的就是把它应用于探测器见证。
我们仅把该引理应用于物理行的无平方因子部分。无平方因子行因子与幂次行因子的
分离也用于 \cite{BGL} 的推论 1.4 的
证明；这里我们保留幂次因子的实际范数，因为它同时控制矩与行的数目。

\begin{proposition}[六次大筛法行包络]\label{prop:sextic-row-count}
固定算术数据 $\mathcal A$、有界的物理范围
$d_{\min}\le d\le d_{\max}$，以及引理~\ref{lem:buffered-bins} 与~\ref{lem:detector-dyads} 的高度与损失假设。令
$U=Z^d$ 并沿用这些引理中的记号 $T_1$。设 $\mathcal B$
为任意一组被保留的无六次幂因子元素行 $u$，满足
$q_u\asymp U$，且全部落在同一个分箱 $(i,a)$ 中，其中 $a>51/100$。令
$\delta=2a-1$，从而 $1/50<\delta\le1$。

对每个 $\epsilon>0$，在减小命题~\ref{prop:detector-witness} 中的预备损失之后，对所有充分大的 $Z$，
\begin{equation}\label{eq:sextic-row-count}
 \#\mathcal B\ll_{\mathcal A,\epsilon}
 U^{R(\delta)+\epsilon}(1+T_1)^{A_{\mathcal A}},
 \qquad
 R(\delta)=\min\left\{1,
       \max\left(1-\frac\delta2,\frac43-\delta\right)\right\}.
\end{equation}
有限的阶 $A_{\mathcal A}$ 关于物理二进区间、分箱以及所有移动的行都是一致的。
它可以在 $T_1=Z^\tau$ 中的正指数被选定之前就固定下来。该估计对
每一组固定的被保留外部参数以及原本的行限制都成立。隐含常数
可以依赖于固定的比较常数以及有限多个共享剖面半范数。逐行的剖面选择
恰为命题~\ref{prop:detector-witness} 与
引理~\ref{lem:smooth-calculus} 所允许的那些；不允许任意的依赖于行的
系数。

下限分箱 $a=51/100$ 被排除在见证断言之外。对该
分箱，初等界 $\#\mathcal B\ll_{\mathcal A}U$ 适用；它与数值
$R(1/50)=1$ 相符。
\end{proposition}

\begin{proof}
固定一个小的预备损失 $\epsilon_0>0$。对 $\mathcal B$ 中的每一行以 $t=1$ 应用
命题~\ref{prop:detector-witness}。按其表示标签
$\psi_{u,\nu,\varsigma}(n)=\nu(n)\chi_n(u)^{\varsigma}$ 及其多项式的二进长度对进行
细分。表示标签只有固定有限多个，
而数对则有 $O_{\mathcal A}((\log U)^2)$ 个。在其中一个这样的
细分内，逆长度 $D=U^r$ 与数对 $(\nu,\varsigma)$
都是固定的。在此细分中记
\[
 M_u(D;W)=D^{-1/2}\sum_n\mu(n)\nu(n)\chi_n(u)^{\varsigma}W(q_n/D)
\]
为该表示的逆多项式。见证给出
\begin{equation}\label{eq:sextic-witness-spike}
 \frac12-O(\epsilon_0)\le r\le1+O(1/\log U),
 \qquad |M_u(D;W_u)|^2\gg_{\mathcal A,\epsilon_0}
       U^{\delta r-\epsilon_0}.
\end{equation}
记号 $W_u$ 保留了被允许的逐行范数剖面参数。
这里的多项式是所展示表示的原本多项式，
而不是其本原诱导特征的多项式。我们首先在这些参数固定时对它取
界，然后处理它们的逐行选择。

对每个无六次幂因子的行，将其理想唯一地分解为
\[
 (u)=\mathfrak s\mathfrak w\mathfrak k,
 \quad
 \mathfrak s=\prod_{p\in S}p^{v_p(u)},\quad
 \mathfrak w=\prod_{\substack{p\notin S\\2\le v_p(u)\le5}}p^{v_p(u)},
 \quad
 \mathfrak k=\prod_{\substack{p\notin S\\v_p(u)=1}}p.
\]
于是 $\mathfrak k$ 无平方因子，$(\mathfrak k,\mathfrak w)=1$，且
$\mathfrak w$ 是幂次的，即其每个正的素
赋值至少为二。$\mathfrak s$ 只有有限多种可能，
因为它在固定集合 $S$ 上的赋值落于 $\{0,\ldots,5\}$。
为这些理想各选取一个生成元 $s_{\mathfrak s}$。以 $w,k$ 记
$\mathfrak w,\mathfrak k$ 的主生成元，则存在唯一的单位 $\upsilon$ 使得
\[
 u=\upsilon s_{\mathfrak s}wk.
\]
每当取行和时，我们都固定 $\mathfrak s$ 与 $\upsilon$。

插入二进区间 $q_{\mathfrak w}\asymp V$，其中把单位理想
包含在有界区间 $V=1$ 内。只可能有 $1\le V\ll_{\mathcal A}U$，
且这样的区间有 $O_{\mathcal A}(\log U)$ 个。原本的圆环
$q_u\asymp U$ 给出实际的剩余圆环
\begin{equation}\label{eq:sextic-actual-row-annulus}
 q_{\mathfrak k}=\frac{q_u}{q_{\mathfrak s}q_{\mathfrak w}}
       \asymp_{\mathcal A}\frac UV.
\end{equation}
在构成和式时我们保留这一限制。只有在应用一个正的行界时，
它才会被放大为球 $q_k\ll_{\mathcal A}U/V$。
以 $\mathcal U_V$ 记当前见证细分中落在这一幂次部分
二进区间内的行，并保留全部原有的物理限制。

范数不超过 $V$ 的幂次理想的数目为 $O(V^{1/2})$。
事实上，每个幂次理想都有唯一的表达式
$\mathfrak v^2\mathfrak y^3$，其中 $\mathfrak y$ 无平方因子；这两个
因子不必互素。由算术预备得到的理想计数给出
\[
 \#\{\mathfrak w:q_{\mathfrak w}\le V,\ \mathfrak w\ {\rm powerful}\}
 \ll V^{1/2}\sum_{\mathfrak y}q_{\mathfrak y}^{-3/2}
 \ll V^{1/2}.
\]
最后一个级数由同样的理想计数收敛。对于其所在二进区间内每一个固定的
$\mathfrak w$，包含
方程~\eqref{eq:sextic-actual-row-annulus} 的球有 $O_{\mathcal A}(U/V)$
个理想。非空的圆环使 $U/V$ 被一个固定正
常数从下界住，这也涵盖了单位剩余理想。由于 $(u)$ 的
分解及其单位都是唯一的，这一幂次部分二进区间内的行的数目满足
\begin{equation}\label{eq:sextic-stratum-cardinality}
 \#\mathcal U_V\ll_{\mathcal A}U V^{-1/2}.
\end{equation}
在这个上界计数中丢弃任何定义 $\mathcal U_V$ 的限制
只会增加行数。

对固定的 $\mathfrak w,\mathfrak s,\upsilon$，令
$v=\upsilon s_{\mathfrak s}w$。分子的乘性给出
精确恒等式
\begin{equation}\label{eq:sextic-frozen-mask}
 \chi_n(u)^{\varsigma}
   =\chi_n(v)^{\varsigma}\chi_n(k)^{\varsigma}.
\end{equation}
它包含了零延拓：若 $n$ 与任一因子有公共因子，则两端
都为零，否则它就是通常的乘性。特别地，
对固定的剖面 $W$，逆多项式是
引理~\ref{lem:sextic-large-sieve} 中的行和，其系数为
\[
 c_n=D^{-1/2}\mu(n)\nu(n)\chi_n(v)^{\varsigma}W(q_n/D).
\]
因子 $\mu(n)$ 使各列保持无平方因子。因子
$\chi_n(v)^{\varsigma}$ 保留了 $(n,v)>1$ 处的每一个原本的零点，且关于
剩余行 $k$ 是固定的。唯一剩下的同时依赖于
$n$ 与 $k$ 的零点是所展示的筛核自身的零点。
限制 $(k,w)=1$，以及对固定外部参数从物理
行集继承来的任何限制，在本应用中都是固定的行限制。
我们不扩大 $S$ 使其包含 $w$ 的素因子。
因此筛中的常数与 $w$ 无关，尽管
其固定的系数掩码可能有很大的范数。

公共的圆环支集与理想计数给出，关于固定的剖面参数一致地，
\[
 \sum_n|c_n|^2
 \le D^{-1}\sum_{q_n\asymp D}|W(q_n/D)|^2\ll_{\mathcal A}1.
\]
同样的陈述对下面用到的有限多个剖面导数也成立，
并带有相应的固定半范数。在把正剩余
行和从其实际圆环放大到 $q_k\ll_{\mathcal A}U/V$ 之后应用
引理~\ref{lem:sextic-large-sieve}。一个固定的
放大使得行界在每个非空的商区间上至少为一，并
包含圆环状的列支集。把结果对
$\mathfrak w$ 的 $O(V^{1/2})$ 种选择与
$\mathfrak s,\upsilon$ 的有限多种选择求和，得到
\begin{equation}\label{eq:sextic-stratum-fixed-moment}
 \sum_{u\in\mathcal U_V}|M_u(D;W)|^2
 \ll_{\mathcal A,\epsilon_1}(UD)^{\epsilon_1}V^{1/2}
 \left\{\frac UV+D+\left(\frac{UD}{V}\right)^{2/3}\right\}.
\end{equation}
这一推导中的所有系数序列在应用筛法的那个行
和内部都是固定的。它们对 $\mathfrak w$ 的不同冻结取值可以不同，
而这些取值只在那些正界之后才被求和。

现在我们恢复方程~\eqref{eq:sextic-witness-spike} 中 $W_u$ 的逐行选择。
这里探测器以 $t=1$ 使用，故 $D_*=U$。在表示与二进对固定
之后，其逆剖面具有形式
\[
 W_{\rm dyad}(x)V_{\le}(Dx/U)x^{-\sigma-i\omega},
 \qquad \omega=\gamma-\nu_F,
\]
其中 $W_{\rm dyad}$ 是固定的光滑二进截断，$\nu_F$ 是
来自探测器的 Fourier 频率。
因子 $V_{\le}(Dx/U)$ 对行和是公共的；变化的参数
是 $\sigma\in[51/100,1]$ 与 $|\omega|\le(3i+1)T_1$。
对任一参数求导会引入一个 $\log x$ 的幂，它在
固定圆环上有界。用单位盒子覆盖这两个参数区间。这样的盒子至多有
$1+T_1$ 的一个固定幂那么多。引理~\ref{lem:smooth-calculus} 中的参数 Sobolev
不等式用剖面导数平方的积分的有限
和界住一个盒子上的上确界。在这些积分之前先对行
求和。在每个固定的参数取值处，求导只改变
固定的剖面系数，引入所允许的对数
权重或剖面导数。因此方程~\eqref{eq:sextic-stratum-fixed-moment}
适用于每个这样的积分。这就证明了
\begin{equation}\label{eq:sextic-stratum-moment}
 \sum_{u\in\mathcal U_V}|M_u(D;W_u)|^2
 \ll_{\mathcal A,\epsilon_1}(UD)^{\epsilon_1}(1+T_1)^{A_{\mathcal A}}
 V^{1/2}\left\{\frac UV+D+
       \left(\frac{UD}{V}\right)^{2/3}\right\}.
\end{equation}
导数的阶以及由此的 $A_{\mathcal A}$ 在选取
$\tau$ 之前就已固定。高度区间与额外频率的
总容许量恰为共享探测器估计中的那些；Sobolev
步骤既不赋予新的频率容许量，也不改变围道。
它不允许对每个 $k$ 分别选取系数。

剩下的是把这个矩与独立的计数
方程~\eqref{eq:sextic-stratum-cardinality} 结合起来。记
\[
 V=U^o,\qquad D=U^r,\qquad
 e_o(r)=\frac o2+\max\left\{1-o,r,\frac{2(1-o+r)}3\right\}.
\]
这些是实际二进区间的参数。非空的二进区间有
$0\le o\le1+O_{\mathcal A}(1/\log U)$；单位二进区间有 $o=0$。
由于见证长度保持有界，方程~\eqref{eq:sextic-stratum-moment} 中 $UD$ 的任意小幂
可以写成 $U$ 的任意小幂。用
方程~\eqref{eq:sextic-witness-spike} 中的尖峰去除该方程，
并同时使用
方程~\eqref{eq:sextic-stratum-cardinality}，对这一细分给出
\begin{equation}\label{eq:sextic-stratum-envelope}
 \#\mathcal U_V
 \ll U^{\min\{1-o/2,e_o(r)-\delta r\}+O(\epsilon_0)+\epsilon_1}
       (1+T_1)^{A_{\mathcal A}},
\end{equation}
在必要时把记作 $\epsilon_1$ 的筛损失减小之后。

我们首先在
$0\le o\le1$ 与 $1/2\le r\le1$ 上计算此表达式中最大的指数。令 $x=1-o$ 且 $A=(1+x)/2$。则
\[
 e_o(r)-\delta r
 =\max\left\{
 A-\delta r,\quad
 \frac{1-x}{2}+(1-\delta)r,\quad
 \frac12+\frac x6+\left(\frac23-\delta\right)r
 \right\}.
\]
$A$ 与三个实数的最大值的最小值，等于 $A$ 分别与这三个数的最小值的最大值。第一个这样的最小值是
$A-\delta r\le1-\delta/2$。对固定的 $r$，第二个是
$x$ 的一个递增与一个递减仿射函数的最小值。它们
在 $x=(1-\delta)r\in[0,1]$ 处相交，故其关于 $x$ 的最大值为
\[
 \frac{1+(1-\delta)r}{2}\le1-\frac\delta2.
\]
第三个最小值中的两个函数都随 $x$ 增大。因此它关于
$x$ 的最大值为
\[
 \min\left\{1,\frac23+\left(\frac23-\delta\right)r\right\}.
\]
这个关于 $r$ 的仿射表达式在两个端点处取值 $1-\delta/2$ 与
$4/3-\delta$。由此推出，且等号在其中一个端点处取得，
\begin{equation}\label{eq:sextic-envelope-optimization}
 \max_{\substack{0\le o\le1\\1/2\le r\le1}}
 \min\{1-o/2,e_o(r)-\delta r\}
 =\min\left\{1,\max\left(1-\frac\delta2,\frac43-\delta\right)\right\}.
\end{equation}
这一计算关于 $0\le\delta\le1$ 是一致成立的。

刚刚用到的函数是仿射函数的最大值与最小值，这些
仿射函数的斜率在这些紧区间的固定邻域上一致有界。
把实际的 $o$ 换成它在 $[0,1]$ 中的最近点，把
实际的见证 $r$ 换成它在 $[1/2,1]$ 中的最近点，由
方程~\eqref{eq:sextic-witness-spike} 与实际二进界，至多使
指数改变 $O_{\mathcal A}(\epsilon_0+1/\log U)$。
因此这一替换只在最后的优化中进行，而不在
行和或其筛法应用中进行。就所指定的 $\epsilon$ 而言，选取 $\epsilon_0$ 与预备
筛损失足够小。$O((\log U)^2)$ 个见证对与 $O(\log U)$ 个幂次部分二进区间只需
再付出一个任意小的幂。对
方程~\eqref{eq:sextic-stratum-envelope} 求和即证得
方程~\eqref{eq:sextic-row-count}。

最后，分箱 $a=51/100$ 不必包含实际的零点，因此那里
不使用命题~\ref{prop:detector-witness}。对其物理圆环中的
所有元素计数得到 $O_{\mathcal A}(U)$ 行。由于
$1/50<1/3$，所展示的 $R$ 的公式给出 $R(1/50)=1$，与断言一致。
\end{proof}

为在高估计中使用，行包络具有如下显式形式
\[
 R(\delta)=
 \begin{cases}
  1,&0\le\delta\le1/3,\\
  4/3-\delta,&1/3\le\delta\le2/3,\\
  1-\delta/2,&2/3\le\delta\le1.
 \end{cases}
\]
论证只用到逆见证。在这个行计数中，不需要关于
与普通见证之乘积的任何估计。

\section{高展开的解析估计}
\label{sec:shared-analytic-estimates}

行包络现在进入基础探针的 Poisson 表示。
对于第一阶段的应应用，保留假设
$\Delta_1=\beta_*-11/12>0$ 并令
$X=Y=Z^{1/2}$。以 $s$ 记
节~\ref{sec:local-euler} 中称为 $x$ 的第一个 Mellin 变量。方程~\eqref{eq:probe-poisson-identity}
与~\eqref{eq:probe-high} 给出待估计的实际积分：
\[
\begin{split}
 I_\eta(Z^{1/2},Z^{1/2},Z)
 =\frac1{(2\pi i)^3}\int_{(3)}\int_{(3)}\int_{(2)}
 &\mathcal W(Z^{1/2},Z^{1/2},Z;s,w,z)\\
 &\cdot\sum_u^{(6)}q_u^{-z}\overline{\xi(u)}
 \frac{\zeta_F^S(6z)L^S(w,\chi_\bullet(u))}
      {L^S(s,\eta\overline{\chi_\bullet(u)})}
 \mathcal H_{\eta,u}(s,w,z)\,dz\,dw\,ds.
\end{split}
\]
这里
\[
 \mathcal W(X,Y,Z;s,w,z)=X^{1/2-z}Z^{s+z-1}Y^{w-1}
 e^{(s+z-1)^2}M(z)\widehat W_1(w).
\]
各行是非零的无六次幂因子元素，带有其单位因子
以及物理限制 $(u,S)=1$。两种特征表示都保留
它们在整除 $u$ 的素数处的原本零点。基础探针的逐系数校准
已包含在这一恒等式中。

行 $u=1$ 包含目标函数的倒数。对于中间行范数，我们将固定一个带缓冲的分箱，把它的积分移到
\[
 \Re s=a+16e,\qquad \Re w=1-a-6e,\qquad \Re z=17/50,
\]
并在那里应用包络 $R(2a-1)$。除小的实部损失与一个固定的高度因子外，
该围道上的分子对行范数 $q_u\asymp U$ 付出 $U^{a-1/2}$ 的代价；行包络控制
这样的项出现多少。我们首先辨识主信号，然后实施
这一移动围道及其具体的第一阶段估计。主留数
与外层行范数的直接界完成了比较。

\subsection{主行与正规化}

我们辨识其分子为主的行。含有一个 $S$ 之外素因子的物理行，
由 \hyperref[par:physical-row-ramification]{节~\ref*{sec:detector} 中的分歧论证}可知其分子非主。
由于物理掩码强制 $(u,S)=1$，
其余的行都是单位。若单位 $u\ne1$ 在 $F$ 中有六次根，
则该根在每个素数处的赋值都为零，从而是 $F$ 的一个
单位；而 $F$ 的每个单位的六次幂都为一。于是
$F(u^{1/6})/F$ 是非平凡的。它是有限 Galois 扩张，因为 $F$ 包含
六次单位根。对非单位 Frobenius 类应用 Chebotarev 定理
可给出一个固定集合 $S$ 之外的素数，使得相应的
六次特征非平凡 \cite[定理 1.1]{TZ}。由
引理~\ref{lem:fixed-numerator-ray} 中的 Kummer 描述，这就是
$\chi_\bullet(u)$。因此 $u=1$ 是唯一的主分子行。
附着于一个有界单位行的分母可能是主的；它的
倒数在一处有零点，将在下面保持在一条整体直线上。

定义
\begin{equation}
 H_\eta(s)=\mathcal H_{\eta,1}(s,1,1/6),\qquad
 c_S=\widehat W_1(1)M(1/6)
       \left(\operatorname*{Res}_{v=1}\zeta_F^S(v)\right)^2/6.
 \label{eq:shared-principal-data}
\end{equation}
引理~\ref{lem:local-euler} 使 $H_\eta$ 在 $\Re s>7/8$ 上全纯。
其一致估计 $H_\eta=1+O(P_0^{-4/5})$ 允许在目标之前
选取 $P_0$，因为局部界关于目标单位
相位是一致的。固定 $P_0$ 足够大，使得
\begin{equation}
 \sup_{\Re s>7/8}|H_\eta(s)-1|\le\frac12.
 \label{eq:shared-principal-nonvanishing}
\end{equation}
用这些素数扩大 $S$ 不改变固定的射线类群 $T$，
因为局部恒等式中的校准是逐系数成立的。被
排除的集合以及由此的特定函数 $H_\eta$ 可以在不同
应用之间不同，但在一个应用内部，它们恰为
所使用的那一个精确高表示的那些。

常数 $c_S$ 为正。事实上 $\widehat W_1(1)>0$，因为
$W_1$ 非负且非零，而由
方程~\eqref{eq:radial-mellin-positive} 有 $M(1/6)>0$。$\zeta_F^S$ 的
单极点具有正的留数：删除把 $\zeta_F$ 的正留数
乘以 $\prod_{p\in S}(1-q_p^{-1})>0$。

对第一阶段，令
\[
 C_{\mathrm I}(s)=s-\frac23,\qquad
 J_{\mathrm I,\eta}(Z)=\frac{I_\eta(Z^{1/2},Z^{1/2},Z)}{c_S}.
\]
设 $f_{\mathrm I,\eta}$ 为
方程~\eqref{eq:common-residue} 中的积分，其中 $C=C_{\mathrm I}$、
$\mathcal S=S$，且函数 $H_\eta$ 如上定义。$w=1$ 与 $z=1/6$ 处的两个标量极点
将给出 $c_S f_{\mathrm I,\eta}$；我们剩余的
目标是为其他每一项取得一个公共的幂次节省，相对于
$Z^{C_{\mathrm I}(\beta_*)}$ 来度量。

\subsection{移动围道所用的精确恒等式}

围道证明使用所示积分中完整的全纯修正。为把
该证明复用于补偿探针，我们记下它所要求的
恒等式与界。在当前的应应用中，修正
是 $\mathcal H_{\eta,u}$，与 $Z$ 无关，而下面的参数
是 $l_x=l_y=h=1/2$ 与 $\ell=0$。

方程~\eqref{eq:probe-high} 中的标量商是高展开中
同时决定移动围道与主留数的那一部分。我们保持该商
固定，并明确陈述围绕它可以改变什么。令
\[
 \begin{split}
 \mathcal D_1(\epsilon_0)=\{(s,w,z):{}&\Re s\ge51/100,\quad
       \Re z\ge17/50,\quad \Re w\ge-1/100,\\
       &\Re(s+w)\ge1+\epsilon_0\},\qquad \epsilon_0>0,\\
 \mathcal D_2=\{(s,w,z):{}&\Re s\ge7/8,\quad
       \Re z\ge33/200,\quad \Re w\ge19/20\}.
 \end{split}
\]
这些是引理~\ref{lem:local-euler} 中的两个区域。下面所谓的实盒子，
指的是 $(\Re s,\Re w,\Re z)$ 的 $\mathbb R^3$ 的一个紧矩形子集。
